Time-series foundation models (TSFMs) forecast unseen series without any training, and their forecast error has been observed to spike when a signal changes. We test whether this makes them zero-shot change point detectors. We propose a forecast-residual detector that scores each time step by its one-step-ahead forecast surprise relative to the model's own predictive interval, built on Chronos-2 and TimesFM-2.5. We compare it with a split-window Chronos-2 embedding detector and with PELT, Binary Segmentation and BOCPD. Evaluation follows the TCPDBench protocol (multi-annotator F1 and segmentation covering, in default and oracle settings) on the 36 real-world series of the Turing Change Point Dataset and the 34 labelled files of the SKAB industrial testbed, and our baselines reproduce TCPDBench's published scores. With parameters fixed in advance, the Chronos-2 residual detector reaches a mean F1 of 0.668 on TCPD. Only Binary Segmentation scores higher (0.686), and the difference is not significant (Wilcoxon test, Holm-corrected $p=0.74$). The detector is strongest on short and medium-length series and weak on series of 512 points or more. When every method is tuned per series, the classical methods are significantly better on both datasets. On SKAB, no detector beats reporting no change points under F1. Zero-shot TSFMs are therefore a usable training-free change point detector, but not a replacement for tuned classical methods.
